\section{Рекомендательный слой и оценка качества}

\subsection{Переход от классификации к ранжированному списку кандидатов}

После выбора основной предиктивной модели следующий шаг состоит в переходе от обычной многоклассовой классификации к рекомендательной постановке. В классификационном режиме модель возвращает наиболее вероятный класс следующей культуры. Однако в прикладном сценарии такой формат не всегда является достаточным. Для специалиста полезнее видеть не один жёсткий ответ, а несколько наиболее правдоподобных кандидатов, которые можно дополнительно интерпретировать с учётом предметного контекста.

По этой причине выход модели рассматривается как список из \(k\) культур-кандидатов.

\[
P(y \mid x) = (p_1, p_2, \ldots, p_m)
\quad \rightarrow \quad
R_k(x) = (c_{(1)}, c_{(2)}, \ldots, c_{(k)}),
\]

где \(c_{(1)}, c_{(2)}, \ldots, c_{(k)}\) --- культуры, соответствующие наибольшим модельным вероятностям.

Такая интерпретация позволяет использовать результат модели не только как единичный прогноз, но и как ранжированный
набор альтернатив для поддержки принятия решений.

\subsection{Метрики рекомендательной постановки}

Для оценки качества в рекомендательном режиме в работе используются три основные метрики:

\begin{itemize}
\tightlist
\item
  \texttt{Accuracy@1}
\item
  \texttt{Hit@3}
\item
  \texttt{Hit@5}
\end{itemize}

Показатель \texttt{Accuracy@1} соответствует точности первой рекомендации и показывает, в какой доле случаев наиболее вероятная культура совпадает с фактической следующей культурой. Показатель \texttt{Hit@3} показывает долю наблюдений, в которых истинная культура попадает в первые три позиции ранжированного списка. Аналогично \texttt{Hit@5} показывает долю наблюдений, в которых правильный класс присутствует в первых пяти рекомендациях.

Содержательно эти метрики позволяют перейти от вопроса \emph{«насколько часто модель угадывает один правильный класс?»} к более прикладному вопросу \emph{«насколько часто правильная культура попадает в короткий список правдоподобных кандидатов?»}

\subsection{Сравнение CatBoost и Markov-1 в рекомендательном режиме}

Для проверки того, насколько полезно использование более длинной истории и контекстных признаков, основная модель CatBoost сравнивается с базовой моделью Markov-1 в той же рекомендательной постановке. Напомним, что Markov-1 использует только последнюю культуру в истории и оценивает переходы вида:

\[
P(c_t \mid c_{t-1})
\]

где \(c_{t-1}\) обозначает последнюю культуру в истории поля, а \(c_t\) --- целевую культуру следующего года.

Таким образом, Markov-1 представляет собой не наивную эвристику, а вероятностную модель переходов между соседними культурами. В отличие от неё CatBoost использует не только последнюю культуру, но и трёхлетнюю историю вместе с контекстными признаками. Сравнение этих подходов показывает вклад более богатого признакового описания в качество рекомендаций.

Итоговые значения для CatBoost и Markov-1 на валидационной и тестовой выборках приведены в таблице~\ref{tab:recommendation-metrics}.

\begin{center}
\small
\begin{longtable}{p{0.24\textwidth} p{0.17\textwidth} r r r}
\caption{Сравнение CatBoost и Markov-1 в рекомендательной постановке}
\label{tab:recommendation-metrics}\\
\toprule
Модель & Выборка & Accuracy@1 & Hit@3 & Hit@5 \\
\midrule
\endfirsthead

\toprule
Модель & Выборка & Accuracy@1 & Hit@3 & Hit@5 \\
\midrule
\endhead

\bottomrule
\endlastfoot

CatBoost & validation & 0.699281 & 0.951190 & 0.990794 \\
Markov-1 & validation & 0.557140 & 0.855302 & 0.931781 \\
CatBoost & test & 0.699006 & 0.951236 & 0.990819 \\
Markov-1 & test & 0.557219 & 0.855303 & 0.931717 \\

\end{longtable}
\footnotesize
Источник значений: \path{artifacts/results/recommendation/recommendation_model_comparison_metrics.csv}.
\normalsize
\end{center}


Полученные значения показывают, что даже если ограничиться одной рекомендацией, CatBoost правильно определяет следующую культуру примерно в 70\% случаев. Однако более важным является то, что при переходе к постановке с формированием короткого списка кандидатов покрытие резко возрастает: в первые три позиции правильная культура попадает примерно в 95\% случаев, а в первые пять позиций --- почти в 99\% случаев.

С прикладной точки зрения это означает, что CatBoost выбран не только по \texttt{Accuracy@1}, но и по способности
формировать качественный список из \(k\) кандидатов для поддержки принятия решений.

Сравнение с Markov-1 показывает устойчивое преимущество основной модели CatBoost по всем ключевым метрикам. Это
означает, что использование трёхлетней истории и контекстных признаков добавляет полезную информацию по сравнению с
моделью, опирающейся только на последнюю культуру. Иными словами, следующая культура определяется не только ближайшим
предыдущим состоянием, но и более широким историческим и контекстным окружением.

\subsection{Поклассовый анализ качества рекомендаций}

Общих метрик недостаточно для полного понимания поведения модели, поэтому в работе дополнительно анализируется качество по отдельным целевым классам. Поклассовый анализ основан на сравнении значений \texttt{recall@1} и \texttt{recall@3} для отдельных целевых классов. Соответствующие значения приведены в таблице~\ref{tab:per-class-recommendation-metrics}.

\begin{table}[H]
\centering
\small
\caption{Поклассовое качество рекомендаций основной модели CatBoost на тестовой выборке}
\label{tab:per-class-recommendation-metrics}
\begin{tabular}{lrrp{0.40\textwidth}}
\toprule
Класс & \texttt{recall@1} & \texttt{recall@3} & Комментарий \\
\midrule
\texttt{corn} & 0.7187 & 0.9841 & Устойчивый массовый класс \\
\texttt{cotton} & 0.7659 & 0.9665 & Хорошо выраженный региональный паттерн \\
\texttt{fallow} & 0.4068 & 0.8044 & Заметный выигрыш при переходе к списку кандидатов \\
\texttt{forage\_hay} & 0.8505 & 0.9337 & Высокое качество первой рекомендации \\
\texttt{legumes} & 0.1384 & 0.7331 & Сложный класс для первой рекомендации \\
\texttt{other\_cereals} & 0.2401 & 0.6394 & Более разнородная группа культур \\
\texttt{sorghum} & 0.3061 & 0.8448 & Заметный выигрыш при переходе к списку кандидатов \\
\texttt{soybeans} & 0.7605 & 0.9928 & Устойчивый массовый класс \\
\texttt{wheat} & 0.7304 & 0.9558 & Один из наиболее устойчивых классов \\
\bottomrule
\end{tabular}

\end{table}


Из таблицы видно, что качество рекомендаций различается по классам. Для устойчивых массовых культур, таких как
\texttt{forage\_hay}, \texttt{cotton}, \texttt{soybeans}, \texttt{wheat} и \texttt{corn}, качество первой рекомендации
уже находится на сравнительно высоком уровне. Для более сложных и разнородных классов, включая \texttt{fallow},
\texttt{sorghum}, \texttt{other\_cereals} и \texttt{legumes}, особенно важен переход к списку из нескольких кандидатов:
даже если первая позиция оказывается ошибочной, правильный класс часто сохраняется среди ближайших альтернатив.

Например, для \texttt{legumes} качество первой рекомендации остаётся низким, однако уже в первых трёх позициях
правильный класс попадает существенно чаще. Это показывает, что в сложных случаях полезность системы заключается не
столько в однозначном угадывании первой позиции, сколько в способности включить правдоподобную культуру в короткий
список альтернатив.

Таким образом, поклассовый анализ подтверждает общий вывод главы: качество рекомендаций зависит от устойчивости
культурных паттернов и представленности классов, а рекомендательная постановка особенно полезна для менее устойчивых
классов.

\subsection{Интерпретация результатов}

Полученные результаты показывают, что рекомендательная постановка лучше соответствует прикладному сценарию использования
системы, чем рассмотрение только первой рекомендации. Высокие значения \texttt{Hit@3} и \texttt{Hit@5} означают, что
модель часто сохраняет правильную культуру среди ближайших альтернатив даже в тех случаях, когда первая позиция
оказывается ошибочной.

При этом сравнение с Markov-1 показывает, что качество рекомендаций определяется не только статистикой ближайшего
перехода \(P(c_t \mid c_{t-1})\), но и более широким историческим и пространственно-региональным контекстом. Это
подтверждает роль CatBoost как предиктивного ядра, формирующего основу для ранжированного списка кандидатов.

В результате рекомендательный слой связывает предиктивное ядро системы с последующими компонентами архитектуры:
анализом уверенности и формированием интерпретационного контекста.